\ifdefined\inMainDoc\else
\documentclass[12pt,a4paper]{article}
% ================================================================
%  PRÉAMBULE COMMUN - ANNALES CORRIGÉES
%  Université de Kara - FaST - Licence Fondamentale MPC
%  Design optimisé impression noir et blanc
% ================================================================

% -- ENCODAGE & LANGUE --
\usepackage[utf8]{inputenc}
\usepackage[T1]{fontenc}
\usepackage[french]{babel}
\usepackage{lmodern}
\usepackage{microtype}

% -- MISE EN PAGE --
\usepackage[a4paper, top=2.8cm, bottom=2.5cm, left=2.5cm, right=2.5cm, headheight=14pt]{geometry}
\usepackage{setspace}
\setstretch{1.2}
\usepackage{parskip}
\setlength{\parindent}{1.2em}
\setlength{\parskip}{4pt}

% -- MATHS --
\usepackage{amsmath, amssymb, mathtools, amsthm}
\usepackage{mathrsfs}
\usepackage{dsfont}
\usepackage{bm}
\usepackage{cancel}
\usepackage{textcomp}
% Définition de \oiint (intégrale de surface fermée) sans package supplémentaire
\newcommand{\oiint}{\mathop{\iint\!\!\!\!\!\!\!\bigcirc}\limits}

% -- PHYSIQUE & CHIMIE --
\usepackage{physics}
\usepackage{siunitx}
% Lorsque physics est chargé, siunitx omet \qty. On le restaure ici :
\AtBeginDocument{\RenewCommandCopy\qty\SI}
\sisetup{locale = FR, detect-all, output-decimal-marker={,}}
% Commande de substitution pour les unités posant problème avec pdflatex
% (ohm, degreeCelsius, micro, angstrom, percent utilisent des caractères non-ASCII)
\newcommand{\qtyohm}[1]{\ensuremath{\num{#1}\,\Omega}}
\newcommand{\qtydegc}[1]{\ensuremath{\num{#1}\,{}^\circ\text{C}}}
\newcommand{\qtyangstrom}[1]{\ensuremath{\num{#1}\,\text{\AA}}}
\newcommand{\qtypercent}[1]{\ensuremath{\num{#1}\,\%}}
\newcommand{\qtyum}[1]{\ensuremath{\num{#1}\,\mu\text{m}}}
\usepackage[version=4]{mhchem}

% -- GRAPHISMES --
\usepackage{graphicx}
\usepackage{float}
\usepackage{caption}
\usepackage{subcaption}
\usepackage{wrapfig}
\usepackage{tikz}
\usetikzlibrary{calc, arrows.meta, patterns, shapes.geometric}
\usepackage{pgfplots}
\pgfplotsset{compat=1.18}

% -- TABLEAUX --
\usepackage{booktabs}
\usepackage{array}
\usepackage{tabularx}
\usepackage{multirow}

% -- LISTES --
\usepackage{enumitem}
\setlist[enumerate]{topsep=3pt, itemsep=2pt, parsep=0pt}
\setlist[itemize]{topsep=3pt, itemsep=2pt, parsep=0pt, label={.}}

% -- BOÎTES (noir et blanc) --
\usepackage{mdframed}
\usepackage{tcolorbox}
\tcbuselibrary{skins, breakable, theorems}

% Boîte EXERCICE : encadré avec filet épais, fond blanc
\newmdenv[
  linewidth=1.2pt,
  linecolor=black,
  roundcorner=3pt,
  skipabove=10pt,
  skipbelow=6pt,
  innertopmargin=8pt,
  innerbottommargin=8pt,
  innerleftmargin=10pt,
  innerrightmargin=10pt,
  backgroundcolor=white,
  frametitle={\bfseries\small Exercice},
  frametitlealignment=\raggedright,
  frametitleaboveskip=4pt,
  frametitlebelowskip=4pt,
]{boiteexercice}

% Boîte CORRIGÉ : fond gris très clair + filet fin
\newmdenv[
  linewidth=0.6pt,
  linecolor=black,
  leftline=true,
  rightline=false,
  topline=false,
  bottomline=false,
  linecolor=black,
  innerleftmargin=12pt,
  innerrightmargin=8pt,
  innertopmargin=6pt,
  innerbottommargin=6pt,
  backgroundcolor=gray!8,
  skipabove=4pt,
  skipbelow=8pt,
]{boitecorrige}

% Boîte RÉSUMÉ DE COURS : double encadré
\newmdenv[
  linewidth=0.8pt,
  linecolor=black,
  roundcorner=2pt,
  backgroundcolor=gray!12,
  innertopmargin=10pt,
  innerbottommargin=10pt,
  innerleftmargin=12pt,
  innerrightmargin=12pt,
  skipabove=12pt,
  skipbelow=8pt,
]{boiteresume}

% Boîte ATTENTION/REMARQUE : barre gauche épaisse
\newmdenv[
  leftline=true,
  rightline=false,
  topline=false,
  bottomline=false,
  linewidth=3pt,
  linecolor=black,
  backgroundcolor=gray!5,
  innerleftmargin=12pt,
  innerrightmargin=8pt,
  innertopmargin=5pt,
  innerbottommargin=5pt,
  skipabove=6pt,
  skipbelow=6pt,
]{boiteremarque}

% -- DIFFICULTÉ DES EXERCICES --
\newcommand{\facile}{$\star$}
\newcommand{\moyen}{$\star\!\star$}
\newcommand{\difficile}{$\star\!\star\!\star$}
\newcommand{\niveau}[1]{\hfill\textbf{#1}\hspace{2pt}}

% -- EN-TÊTES ET PIEDS DE PAGE --
\usepackage{fancyhdr}
\pagestyle{fancy}
\fancyhf{}
\renewcommand{\headrulewidth}{0.5pt}
\renewcommand{\footrulewidth}{0.3pt}

% -- HYPERLIENS --
\usepackage[hidelinks]{hyperref}
\hypersetup{
    colorlinks=false,
    pdftitle={Annale Corrigée - Licence MPC - Université de Kara}
}

% -- TITRES --
\usepackage{titlesec}
\titleformat{\section}{\Large\bfseries}{\thesection.}{0.8em}{}[\vspace{2pt}\hrule\vspace{4pt}]
\titleformat{\subsection}{\large\bfseries}{\thesubsection.}{0.7em}{}
\titleformat{\subsubsection}{\normalsize\bfseries}{\thesubsubsection.}{0.6em}{}

% -- THÉORÈMES --
\theoremstyle{definition}
\newtheorem{defn}{Définition}[section]
\newtheorem{prop}{Propriété}[section]
\newtheorem{thm}{Théorème}[section]
\newtheorem{corol}{Corollaire}[section]
\newtheorem{exemple}{Exemple}[section]

% -- COMMANDES MATHÉMATIQUES UTILES --
\newcommand{\vect}[1]{\overrightarrow{#1}}
\newcommand{\R}{\mathbb{R}}
\newcommand{\N}{\mathbb{N}}
\newcommand{\Z}{\mathbb{Z}}
\newcommand{\C}{\mathbb{C}}
\renewcommand{\d}{\mathrm{d}}
\newcommand{\deriv}[2]{\dfrac{\d #1}{\d #2}}
\newcommand{\dpart}[2]{\dfrac{\partial #1}{\partial #2}}
\newcommand{\rot}{\overrightarrow{\mathrm{rot}}\,}
\renewcommand{\div}{\mathrm{div}\,}
\renewcommand{\grad}{\overrightarrow{\mathrm{grad}}\,}
\newcommand{\e}{\mathrm{e}}
\newcommand{\ii}{\mathrm{i}}
\renewcommand{\Tr}{\mathrm{Tr}}

% Numérotation des équations par section
\numberwithin{equation}{section}

% -- RÉSULTATS ENCADRÉS --
\newcommand{\res}[1]{\[\boxed{#1}\]}

% -- REMARQUE INLINE --
\newcommand{\rmq}[1]{\begin{boiteremarque}\textbf{Remarque.} #1\end{boiteremarque}}

% -- MISE EN VALEUR DU RÉSULTAT FINAL --
\newcommand{\reponse}[1]{%
  \begin{center}
    \fbox{\begin{minipage}{0.92\textwidth}\centering #1\end{minipage}}
  \end{center}%
}


\lhead{\small\textit{Séries de Fourier et Transformée de Laplace - Licence 2}}
\rhead{\small\textit{FaST, Université de Kara}}
\cfoot{\thepage}

\begin{document}
\fi

\begin{center}
  {\LARGE\bfseries SÉRIES DE FOURIER ET TRANSFORMÉE DE LAPLACE}\\[0.3cm]
  {\normalsize Licence 2 - Mathématiques, Physique et Chimie}\\[0.1cm]
  \rule{0.7\textwidth}{0.8pt}
\end{center}

\section*{Résumé du cours}

\begin{boiteresume}

\textbf{1. Série de Fourier.} Toute fonction $f$ périodique de période $T=2\pi/\omega_0$, intégrable sur une période, admet un développement en série de Fourier :
\[
f(t) = \frac{a_0}{2} + \sum_{n=1}^{+\infty}\left[a_n\cos(n\omega_0 t) + b_n\sin(n\omega_0 t)\right]
\]
avec $a_n = \frac{2}{T}\int_{-T/2}^{T/2} f(t)\cos(n\omega_0 t)\,\d t$ et $b_n = \frac{2}{T}\int_{-T/2}^{T/2} f(t)\sin(n\omega_0 t)\,\d t$.

\textbf{2. Parité.} Si $f$ est paire, $b_n=0$ (série en cosinus). Si $f$ est impaire, $a_n=0$ (série en sinus).

\textbf{3. Convergence.} Si $f$ est continue par morceaux sur $\mathbb{R}$, la série converge vers $\frac{f(t^+)+f(t^-)}{2}$ en tout point, et vers $f(t)$ aux points de continuité.

\textbf{4. Égalité de Parseval.} $\frac{1}{T}\int_0^T |f(t)|^2\,\d t = \frac{a_0^2}{4}+\frac{1}{2}\sum_{n=1}^{+\infty}(a_n^2+b_n^2)$.

\textbf{5. Transformée de Laplace.} $\mathcal{L}\{f\}(s) = F(s) = \int_0^{+\infty} f(t)\,\e^{-st}\,\d t$ (converge pour $\text{Re}(s)>\alpha_0$). Propriétés clés :
\begin{itemize}[nosep]
  \item $\mathcal{L}\{f'\} = sF(s)-f(0)$ ; $\mathcal{L}\{f''\} = s^2F(s)-sf(0)-f'(0)$
  \item $\mathcal{L}\{\e^{at}f\} = F(s-a)$ (décalage en fréquence)
  \item $\mathcal{L}\{f*g\} = F(s)\cdot G(s)$ (convolution)
\end{itemize}
Transformées usuelles : $\mathcal{L}\{1\}=\frac{1}{s}$, $\mathcal{L}\{\e^{at}\}=\frac{1}{s-a}$, $\mathcal{L}\{\sin(\omega t)\}=\frac{\omega}{s^2+\omega^2}$, $\mathcal{L}\{\cos(\omega t)\}=\frac{s}{s^2+\omega^2}$.

\end{boiteresume}

\newpage
\section{Séries de Fourier}

\subsection*{Exercice 1 - Série de Fourier d'un signal créneau \niveau{\facile}}

\begin{boiteexercice}
Soit $f$ la fonction $2\pi$-périodique définie sur $[-\pi,\pi[$ par :
\[
f(t) = \begin{cases}1 & \text{si } t\in[0,\pi[\\-1 & \text{si } t\in[-\pi,0[\end{cases}
\]
\begin{enumerate}
  \item Calculer les coefficients de Fourier $a_n$ et $b_n$.
  \item Écrire le développement en série de Fourier de $f$.
  \item Vers quoi converge la série en $t=0$ et en $t=\pi$ ?
  \item Déduire la valeur de la série $1 - \frac{1}{3} + \frac{1}{5} - \frac{1}{7} + \cdots$
\end{enumerate}
\end{boiteexercice}

\begin{boitecorrige}
\textbf{1.} $f$ est impaire, donc $a_n = 0$ pour tout $n\geq0$.

$b_n = \frac{1}{\pi}\int_{-\pi}^{\pi} f(t)\sin(nt)\,\d t = \frac{2}{\pi}\int_0^\pi \sin(nt)\,\d t = \frac{2}{\pi}\left[-\frac{\cos(nt)}{n}\right]_0^\pi = \frac{2}{n\pi}(1-\cos(n\pi)) = \frac{2}{n\pi}(1-(-1)^n)$.

Donc $b_n = \begin{cases}\frac{4}{n\pi} & \text{si } n \text{ impair}\\0 & \text{si } n \text{ pair}\end{cases}$.

\textbf{2.} Développement :
\[
f(t) = \frac{4}{\pi}\left(\sin t + \frac{\sin 3t}{3} + \frac{\sin 5t}{5} + \cdots\right) = \frac{4}{\pi}\sum_{k=0}^{+\infty}\frac{\sin((2k+1)t)}{2k+1}
\]

\textbf{3.} En $t=0$ : $f$ est continue en 0 avec $f(0^+)=1$ et $f(0^-)=-1$, la série converge vers $\frac{1+(-1)}{2}=0$.

En $t=\pi$ : de même, la série converge vers $\frac{-1+1}{2}=0$.

\textbf{4.} En $t=\pi/2$, $f(\pi/2)=1$ et la série converge vers $f(\pi/2)=1$ (point de continuité) :
\[
\frac{4}{\pi}\left(1 - \frac{1}{3} + \frac{1}{5} - \frac{1}{7}+\cdots\right) = 1 \quad\Rightarrow\quad \sum_{k=0}^{+\infty}\frac{(-1)^k}{2k+1} = \frac{\pi}{4}
\]
C'est la formule de Leibniz-Gregory.
\end{boitecorrige}

\subsection*{Exercice 2 - Série de Fourier d'une fonction triangulaire \niveau{\moyen}}

\begin{boiteexercice}
Soit $f$ la fonction $2\pi$-périodique définie sur $(-\pi,\pi]$ par $f(t) = |t|$.
\begin{enumerate}
  \item Calculer les coefficients $a_n$ et $b_n$.
  \item Écrire le développement en série de Fourier.
  \item Déduire la valeur de $\sum_{n=0}^{+\infty}\frac{1}{(2n+1)^2}$.
  \item Appliquer l'égalité de Parseval et en déduire $\sum_{n=1}^{+\infty}\frac{1}{n^2}$.
\end{enumerate}
\end{boiteexercice}

\begin{boitecorrige}
\textbf{1.} $f$ est paire, donc $b_n=0$. $a_0 = \frac{2}{\pi}\int_0^\pi t\,\d t = \frac{2}{\pi}\cdot\frac{\pi^2}{2} = \pi$.

Pour $n\geq1$ : $a_n = \frac{2}{\pi}\int_0^\pi t\cos(nt)\,\d t$. Par parties :
\[
a_n = \frac{2}{\pi}\left[\frac{t\sin(nt)}{n}\right]_0^\pi - \frac{2}{n\pi}\int_0^\pi\sin(nt)\,\d t = 0 + \frac{2}{n^2\pi}[\cos(nt)]_0^\pi = \frac{2((-1)^n-1)}{n^2\pi}
\]
Donc $a_n=0$ si $n$ est pair, et $a_n = -\frac{4}{n^2\pi}$ si $n$ est impair.

\textbf{2.}
\[
f(t) = \frac{\pi}{2} - \frac{4}{\pi}\left(\cos t + \frac{\cos 3t}{9} + \frac{\cos 5t}{25}+\cdots\right) = \frac{\pi}{2} - \frac{4}{\pi}\sum_{k=0}^{+\infty}\frac{\cos((2k+1)t)}{(2k+1)^2}
\]

\textbf{3.} En $t=0$ : $f(0)=0$, donc $0 = \frac{\pi}{2} - \frac{4}{\pi}\sum_{k=0}^{+\infty}\frac{1}{(2k+1)^2}$, d'où $\displaystyle\sum_{k=0}^{+\infty}\frac{1}{(2k+1)^2} = \frac{\pi^2}{8}$.

\textbf{4.} Parseval : $\frac{1}{2\pi}\int_{-\pi}^\pi t^2\,\d t = \frac{a_0^2}{4}+\frac{1}{2}\sum_{n\text{ impair}}\frac{16}{n^4\pi^2}$.

LHS : $\frac{\pi^2}{3}$. On en déduit $\frac{\pi^2}{3} - \frac{\pi^2}{4} = \frac{8}{\pi^2}\sum_{n=0}^{+\infty}\frac{1}{(2n+1)^4}$... En utilisant $\sum_{n=1}^\infty\frac{1}{n^2} = \sum_{n=1}^\infty\frac{1}{(2n)^2}+\sum_{n=0}^\infty\frac{1}{(2n+1)^2} = \frac{1}{4}\sum_{n=1}^\infty\frac{1}{n^2}+\frac{\pi^2}{8}$, on obtient $\frac{3}{4}\sum_{n=1}^\infty\frac{1}{n^2}=\frac{\pi^2}{8}$ d'où :
\[
\sum_{n=1}^{+\infty}\frac{1}{n^2} = \frac{\pi^2}{6}
\]
\end{boitecorrige}

\section{Transformée de Laplace}

\subsection*{Exercice 3 - Calcul de transformées \niveau{\facile}}

\begin{boiteexercice}
Calculer les transformées de Laplace suivantes :
\begin{enumerate}[label=(\alph*)]
  \item $\mathcal{L}\{t^2\}(s)$
  \item $\mathcal{L}\{\e^{-2t}\sin(3t)\}(s)$
  \item $\mathcal{L}\{t\e^{at}\}(s)$
  \item $\mathcal{L}\{\cosh(at)\}(s)$
\end{enumerate}
\end{boiteexercice}

\begin{boitecorrige}
\textbf{(a)} $\mathcal{L}\{t^n\}(s) = n!/s^{n+1}$, donc $\mathcal{L}\{t^2\} = \dfrac{2}{s^3}$.

\textbf{(b)} $\mathcal{L}\{\sin(3t)\} = \frac{3}{s^2+9}$. Par décalage ($a=-2$) : $\mathcal{L}\{\e^{-2t}\sin(3t)\} = \dfrac{3}{(s+2)^2+9}$.

\textbf{(c)} $\mathcal{L}\{t\} = 1/s^2$. Par décalage : $\mathcal{L}\{t\e^{at}\} = \dfrac{1}{(s-a)^2}$.

\textbf{(d)} $\cosh(at)=\frac{\e^{at}+\e^{-at}}{2}$, donc $\mathcal{L}\{\cosh(at)\} = \frac{1}{2}\left(\frac{1}{s-a}+\frac{1}{s+a}\right) = \dfrac{s}{s^2-a^2}$.
\end{boitecorrige}

\subsection*{Exercice 4 - Résolution d'EDO par Laplace \niveau{\moyen}}

\begin{boiteexercice}
Résoudre par la transformée de Laplace :
\begin{enumerate}[label=(\alph*)]
  \item $y''-3y'+2y=0$, avec $y(0)=1$, $y'(0)=0$.
  \item $y''+4y=\sin(2t)$, avec $y(0)=0$, $y'(0)=1$.
\end{enumerate}
\end{boiteexercice}

\begin{boitecorrige}
\textbf{(a)} Notons $Y=\mathcal{L}\{y\}$. On applique $\mathcal{L}$ à l'équation :
\[
s^2Y-s\cdot1-0 - 3(sY-1) + 2Y = 0 \quad\Rightarrow\quad Y(s^2-3s+2) = s-3
\]
\[
Y = \frac{s-3}{(s-1)(s-2)} = \frac{A}{s-1}+\frac{B}{s-2}
\]
$A = \frac{1-3}{1-2}=2$, $B=\frac{2-3}{2-1}=-1$. Par transformée inverse :
\[
\boxed{y(t) = 2\e^t - \e^{2t}}
\]

\textbf{(b)} $s^2Y - s\cdot0 - 1 + 4Y = \frac{2}{s^2+4}$ :
\[
Y(s^2+4) = 1+\frac{2}{s^2+4} \quad\Rightarrow\quad Y = \frac{1}{s^2+4}+\frac{2}{(s^2+4)^2}
\]
$\mathcal{L}^{-1}\left\{\frac{1}{s^2+4}\right\} = \frac{\sin(2t)}{2}$.

$\mathcal{L}^{-1}\left\{\frac{2}{(s^2+4)^2}\right\}$ : par la table, $\mathcal{L}^{-1}\left\{\frac{2\omega}{(s^2+\omega^2)^2}\right\} = \frac{\sin(\omega t)-\omega t\cos(\omega t)}{2\omega^2}$. Avec $\omega=2$ :

$\mathcal{L}^{-1}\left\{\frac{2}{(s^2+4)^2}\right\} = \frac{\sin(2t)-2t\cos(2t)}{8}$.

\[
\boxed{y(t) = \frac{\sin(2t)}{2} + \frac{\sin(2t)-2t\cos(2t)}{8} = \frac{5\sin(2t)}{8} - \frac{t\cos(2t)}{4}}
\]
Le terme $t\cos(2t)$ traduit le phénomène de résonance.
\end{boitecorrige}

\subsection*{Exercice 5 - Application : circuit RLC \niveau{\difficile}}

\begin{boiteexercice}
Un circuit RLC série est décrit par $L\ddot{q} + R\dot{q} + q/C = e(t)$ avec $L=1$\,H, $R=2$\,$\Omega$, $C=1$\,F, $e(t)=U_0[1-\e^{-t}]$ pour $t>0$, et les conditions initiales nulles $q(0)=\dot{q}(0)=0$.
\begin{enumerate}
  \item Écrire l'équation en $q$.
  \item Résoudre par la transformée de Laplace.
  \item Interpréter physiquement la solution.
\end{enumerate}
\end{boiteexercice}

\begin{boitecorrige}
\textbf{1.} $\ddot{q} + 2\dot{q} + q = U_0(1-\e^{-t})$.

\textbf{2.} Avec $Q=\mathcal{L}\{q\}$ et $E(s)=\mathcal{L}\{U_0(1-\e^{-t})\}=U_0\left(\frac{1}{s}-\frac{1}{s+1}\right)=\frac{U_0}{s(s+1)}$ :
\[
(s^2+2s+1)Q = \frac{U_0}{s(s+1)} \quad\Rightarrow\quad Q = \frac{U_0}{s(s+1)^3}
\]
Par décomposition en éléments simples :
\[
Q = U_0\left[\frac{1}{s} - \frac{1}{s+1} - \frac{1}{(s+1)^2} - \frac{1}{2(s+1)^3}\right]
\]
En inversant :
\[
q(t) = U_0\left[1 - \e^{-t} - t\e^{-t} - \frac{t^2}{2}\e^{-t}\right] = U_0\left[1 - \e^{-t}\left(1+t+\frac{t^2}{2}\right)\right]
\]

\textbf{3.} Le système est amorti ($R=2\Omega$, racine double $s=-1$). La charge $q(t)$ tend asymptotiquement vers $U_0$ (régime permanent). Le terme $\e^{-t}(1+t+t^2/2)$ représente le régime transitoire qui s'atténue exponentiellement.
\end{boitecorrige}

\section{Transformée de Fourier}

\subsection*{Exercice 6 - Transformée de Fourier : signal porte et signal exponentiel \niveau{\moyen}}

\begin{boiteexercice}
On rappelle que la transformée de Fourier est définie par $\hat{f}(\nu) = \int_{-\infty}^{+\infty} f(t)\,\e^{-2\pi i\nu t}\,\d t$.
\begin{enumerate}
  \item Calculer la TF du signal porte : $f(t) = \begin{cases}1 & \text{si } |t| \leq T/2\\0 & \text{sinon}\end{cases}$.
  \item Calculer la TF du signal exponentiel causal : $g(t) = \e^{-at}\mathbf{1}_{t\geq0}$, $a > 0$.
  \item Vérifier que $|\hat{g}(\nu)|^2$ représente un spectre lorentzien et interpréter $a$.
\end{enumerate}
\end{boiteexercice}

\begin{boitecorrige}
\textbf{1.} Signal porte :
\[
\hat{f}(\nu) = \int_{-T/2}^{T/2} \e^{-2\pi i\nu t}\,\d t = \left[\frac{\e^{-2\pi i\nu t}}{-2\pi i\nu}\right]_{-T/2}^{T/2} = \frac{\e^{-i\pi\nu T} - \e^{i\pi\nu T}}{-2\pi i\nu} = \frac{2\sin(\pi\nu T)}{2\pi\nu}
\]
\[
\boxed{\hat{f}(\nu) = T\,\mathrm{sinc}(\pi\nu T) = \frac{\sin(\pi\nu T)}{\pi\nu}}
\]
où $\mathrm{sinc}(x) = \sin(x)/x$. La TF d'une porte est un sinus cardinal. Plus la porte est large (grand $T$), plus le spectre est étroit (principe d'incertitude).

\textbf{2.} Signal exponentiel causal :
\[
\hat{g}(\nu) = \int_0^{+\infty}\e^{-at}\e^{-2\pi i\nu t}\,\d t = \int_0^{+\infty}\e^{-(a+2\pi i\nu)t}\,\d t = \left[\frac{\e^{-(a+2\pi i\nu)t}}{-(a+2\pi i\nu)}\right]_0^{+\infty} = \frac{1}{a+2\pi i\nu}
\]
(converge car $a > 0$). Donc :
\[
\boxed{\hat{g}(\nu) = \frac{1}{a+2\pi i\nu}}
\]

\textbf{3.} Le module au carré (densité spectrale de puissance) :
\[
|\hat{g}(\nu)|^2 = \frac{1}{a^2+(2\pi\nu)^2} = \frac{1}{a^2}\cdot\frac{1}{1+(\nu/(\frac{a}{2\pi}))^2}
\]
C'est une lorentzienne centrée en $\nu = 0$ de largeur à mi-hauteur $\Delta\nu = a/\pi$. Le paramètre $a$ est l'inverse du temps de vie : plus $a$ est petit (signal qui dure longtemps), plus le spectre est étroit.
\end{boitecorrige}

\subsection*{Exercice 7 - Propriétés de la transformée de Fourier \niveau{\moyen}}

\begin{boiteexercice}
On note $\hat{f}(\nu) = \mathcal{F}\{f\}(\nu)$ et on rappelle les propriétés :
\begin{itemize}[nosep]
  \item Décalage temporel : $\mathcal{F}\{f(t-t_0)\}(\nu) = \e^{-2\pi i\nu t_0}\hat{f}(\nu)$.
  \item Dérivation : $\mathcal{F}\{f'(t)\}(\nu) = 2\pi i\nu\,\hat{f}(\nu)$.
  \item Homothétie : $\mathcal{F}\{f(at)\}(\nu) = \dfrac{1}{|a|}\hat{f}(\nu/a)$.
\end{itemize}
\begin{enumerate}
  \item On sait que $\mathcal{F}\{\e^{-\pi t^2}\}(\nu) = \e^{-\pi\nu^2}$ (gaussienne auto-transformée). En déduire $\mathcal{F}\{\e^{-\pi t^2/\sigma^2}\}(\nu)$.
  \item En utilisant la propriété de dérivation, calculer $\mathcal{F}\{t\,\e^{-\pi t^2}\}(\nu)$ en dérivant par rapport à $\nu$ l'identité $\hat{f}(\nu) = \e^{-\pi\nu^2}$.
  \item Soit $f(t) = \cos(2\pi f_0 t)\,\e^{-at}$ pour $t\geq0$ et 0 sinon. Exprimer $\hat{f}(\nu)$ en utilisant la propriété de décalage fréquentiel (modulation).
\end{enumerate}
\end{boiteexercice}

\begin{boitecorrige}
\textbf{1.} On utilise la propriété d'homothétie avec $a = 1/\sigma$ :
\[
\mathcal{F}\!\left\{\e^{-\pi t^2/\sigma^2}\right\}(\nu) = \mathcal{F}\!\left\{\e^{-\pi (t/\sigma)^2}\right\}(\nu) = |\sigma|\,\hat{f}(\sigma\nu) = \sigma\,\e^{-\pi\sigma^2\nu^2}
\]

\textbf{2.} En dérivant $\hat{f}(\nu) = \int_{-\infty}^{+\infty}\e^{-\pi t^2}\e^{-2\pi i\nu t}\,\d t$ par rapport à $\nu$ (dérivation sous le signe intégral) :
\[
\frac{\d\hat{f}}{\d\nu} = \int_{-\infty}^{+\infty}\e^{-\pi t^2}(-2\pi it)\e^{-2\pi i\nu t}\,\d t = -2\pi i\,\mathcal{F}\{t\,\e^{-\pi t^2}\}(\nu)
\]
Or $\hat{f}(\nu) = \e^{-\pi\nu^2}$, donc $\dfrac{\d\hat{f}}{\d\nu} = -2\pi\nu\,\e^{-\pi\nu^2}$, d'où :
\[
\mathcal{F}\{t\,\e^{-\pi t^2}\}(\nu) = \frac{-(-2\pi\nu)\e^{-\pi\nu^2}}{2\pi i} = \frac{-i\nu\,\e^{-\pi\nu^2}}{1} = -i\nu\,\e^{-\pi\nu^2}
\]

\textbf{3.} On écrit $\cos(2\pi f_0 t) = \dfrac{\e^{2\pi if_0 t}+\e^{-2\pi if_0 t}}{2}$. Alors :
\[
f(t) = \frac{1}{2}\left(\e^{(2\pi if_0 - a)t} + \e^{-(2\pi if_0+a)t}\right)\mathbf{1}_{t\geq0}
\]
\[
\hat{f}(\nu) = \frac{1}{2}\left(\frac{1}{a-2\pi if_0 + 2\pi i\nu} + \frac{1}{a+2\pi if_0+2\pi i\nu}\right)\cdot(-1)\times(-1)
\]
En utilisant $\hat{g}(\nu) = 1/(a+2\pi i\nu)$ et le décalage en fréquence :
\[
\hat{f}(\nu) = \frac{1}{2}\left[\frac{1}{a+2\pi i(\nu-f_0)} + \frac{1}{a+2\pi i(\nu+f_0)}\right]
\]
Le spectre est une superposition de deux lorentziennes centrées en $\pm f_0$.
\end{boitecorrige}

\subsection*{Exercice 8 - Convolution et théorème de convolution \niveau{\difficile}}

\begin{boiteexercice}
La convolution est définie par $(f*g)(t) = \int_{-\infty}^{+\infty}f(\tau)g(t-\tau)\,\d\tau$. Le théorème de convolution affirme : $\mathcal{F}\{f*g\} = \hat{f}\cdot\hat{g}$.
\begin{enumerate}
  \item Calculer la convolution $h = f*f$ où $f(t) = \mathbf{1}_{[0,1]}(t)$ (indicatrice de $[0,1]$) en discutant les cas.
  \item Vérifier par la transformée de Fourier : calculer $\hat{f}(\nu)^2$ et identifier la TF inverse.
  \item Soit un système LTI (linéaire invariant dans le temps) de réponse impulsionnelle $h(t) = \e^{-t}\mathbf{1}_{t\geq0}$. Calculer sa fonction de transfert $H(\nu) = \hat{h}(\nu)$.
  \item Quelle est la sortie $y(t)$ si l'entrée est $x(t) = \e^{-2t}\mathbf{1}_{t\geq0}$ ? (Utiliser la TF.)
\end{enumerate}
\end{boiteexercice}

\begin{boitecorrige}
\textbf{1.} $h(t) = (f*f)(t) = \int_0^1 f(t-\tau)\,\d\tau = \int_0^1 \mathbf{1}_{[0,1]}(t-\tau)\,\d\tau$.

$t-\tau\in[0,1] \Leftrightarrow t-1\leq\tau\leq t$. On distingue :
\begin{itemize}[nosep]
  \item $t < 0$ ou $t > 2$ : $h(t) = 0$.
  \item $0 \leq t \leq 1$ : $\int_0^t\d\tau = t$.
  \item $1 < t \leq 2$ : $\int_{t-1}^1\d\tau = 2-t$.
\end{itemize}
Donc $h(t) = \begin{cases}t & 0\leq t\leq1\\2-t & 1<t\leq2\\0 & \text{sinon}\end{cases}$ : c'est une fonction chapeau (triangle).

\textbf{2.} $\hat{f}(\nu) = \dfrac{1-\e^{-2\pi i\nu}}{2\pi i\nu}$ (porte de largeur 1). Alors :
\[
\hat{f}(\nu)^2 = \frac{(1-\e^{-2\pi i\nu})^2}{(2\pi i\nu)^2} = \frac{1-2\e^{-2\pi i\nu}+\e^{-4\pi i\nu}}{-4\pi^2\nu^2}
\]
On peut vérifier que la TF inverse donne bien la fonction chapeau.

\textbf{3.} $H(\nu) = \hat{h}(\nu) = \dfrac{1}{1+2\pi i\nu}$.

\textbf{4.} $X(\nu) = \hat{x}(\nu) = \dfrac{1}{2+2\pi i\nu}$. Par le théorème de convolution :
\[
Y(\nu) = H(\nu)\cdot X(\nu) = \frac{1}{(1+2\pi i\nu)(2+2\pi i\nu)}
\]
Décomposition en éléments simples : $Y(\nu) = \dfrac{1}{1+2\pi i\nu} - \dfrac{1}{2+2\pi i\nu}$. Par TF inverse :
\[
y(t) = (\e^{-t} - \e^{-2t})\,\mathbf{1}_{t\geq0}
\]
Ce résultat est cohérent avec la résolution directe de l'équation différentielle $y' + y = x$.
\end{boitecorrige}

\subsection*{Exercice 9 - Parseval et calcul de sommes \niveau{\moyen}}

\begin{boiteexercice}
On rappelle l'égalité de Parseval pour les séries de Fourier : si $f$ a pour développement de Fourier $\dfrac{a_0}{2}+\sum_{n=1}^\infty(a_n\cos(nt)+b_n\sin(nt))$, alors :
\[
\frac{1}{\pi}\int_0^{2\pi}|f(t)|^2\,\d t = \frac{a_0^2}{2} + \sum_{n=1}^{+\infty}(a_n^2+b_n^2)
\]
\begin{enumerate}
  \item En appliquant Parseval à la série de Fourier du signal créneau (exercice 1), calculer $\displaystyle\sum_{n=0}^{+\infty}\frac{1}{(2n+1)^2}$.
  \item En appliquant Parseval à $f(t) = t$ sur $[-\pi,\pi]$ (dont le développement est $f(t)=2\sum_{n=1}^\infty\dfrac{(-1)^{n+1}}{n}\sin(nt)$), calculer $\displaystyle\sum_{n=1}^{+\infty}\frac{1}{n^2}$.
  \item Déduire la valeur de $\displaystyle\sum_{n=1}^{+\infty}\frac{1}{n^4}$ en appliquant Parseval à $f(t)=t^2$ sur $[-\pi,\pi]$.
\end{enumerate}
\end{boiteexercice}

\begin{boitecorrige}
\textbf{1.} Le signal créneau a $a_n = 0$ et $b_n = 4/(n\pi)$ pour $n$ impair, $b_n = 0$ pour $n$ pair. Parseval donne :
\[
\frac{1}{\pi}\int_{-\pi}^\pi 1^2\,\d t = \sum_{k=0}^{+\infty}b_{2k+1}^2 \quad\Rightarrow\quad 2 = \sum_{k=0}^{+\infty}\frac{16}{(2k+1)^2\pi^2}
\]
\[
\boxed{\sum_{k=0}^{+\infty}\frac{1}{(2k+1)^2} = \frac{\pi^2}{8}}
\]

\textbf{2.} Pour $f(t)=t$ : Parseval donne $\dfrac{1}{\pi}\int_{-\pi}^\pi t^2\,\d t = \sum_{n=1}^\infty\frac{4}{n^2}$. LHS $= \dfrac{1}{\pi}\cdot\dfrac{2\pi^3}{3} = \dfrac{2\pi^2}{3}$. Donc :
\[
\frac{2\pi^2}{3} = 4\sum_{n=1}^\infty\frac{1}{n^2} \quad\Rightarrow\quad \boxed{\sum_{n=1}^{+\infty}\frac{1}{n^2} = \frac{\pi^2}{6}}
\]

\textbf{3.} Pour $f(t) = t^2$ sur $[-\pi,\pi]$, le développement
en série de Fourier est :
\[
\dfrac{\pi^2}{3}+4\sum_{n=1}^\infty\dfrac{(-1)^n}{n^2}\cos(nt).
\]
Parseval :
\[
\frac{1}{\pi}\int_{-\pi}^\pi t^4\,\d t = \frac{a_0^2}{2}+\sum_{n=1}^\infty a_n^2
\]
LHS : $\dfrac{1}{\pi}\cdot\dfrac{2\pi^5}{5} = \dfrac{2\pi^4}{5}$. Terme $a_0 = 2\pi^2/3$, $a_n = 4(-1)^n/n^2$.
\[
\frac{2\pi^4}{5} = \frac{1}{2}\!\left(\frac{2\pi^2}{3}\right)^2 + \sum_{n=1}^\infty\frac{16}{n^4} = \frac{2\pi^4}{9} + 16\sum_{n=1}^\infty\frac{1}{n^4}
\]
\[
\sum_{n=1}^{+\infty}\frac{1}{n^4} = \frac{1}{16}\!\left(\frac{2\pi^4}{5} - \frac{2\pi^4}{9}\right) = \frac{2\pi^4}{16}\cdot\frac{9-5}{45} = \frac{2\pi^4}{16}\cdot\frac{4}{45} = \boxed{\frac{\pi^4}{90}}
\]
\end{boitecorrige}

\subsection*{Exercice 10 - Série de Fourier d'une fonction affine par morceaux \niveau{\moyen}}

\begin{boiteexercice}
Soit $f$ la fonction $2\pi$-périodique définie sur $[0, 2\pi[$ par $f(t) = \pi - t$.
\begin{enumerate}
  \item Calculer $a_0$, $a_n$ et $b_n$ (pour $n\geq1$).
  \item Écrire le développement de Fourier de $f$.
  \item Vers quoi converge la série en $t = 0$ ? Interpréter graphiquement (phénomène de Gibbs).
  \item En appliquant l'identité de Parseval, calculer $\displaystyle\sum_{n=1}^{+\infty}\frac{1}{n^2}$.
\end{enumerate}
\end{boiteexercice}

\begin{boitecorrige}
\textbf{1.} $a_0 = \dfrac{1}{\pi}\int_0^{2\pi}(\pi-t)\,\d t = \dfrac{1}{\pi}\left[\pi t - \dfrac{t^2}{2}\right]_0^{2\pi} = \dfrac{1}{\pi}(2\pi^2 - 2\pi^2) = 0$.

Pour $n\geq1$ :
\[
a_n = \frac{1}{\pi}\int_0^{2\pi}(\pi-t)\cos(nt)\,\d t = \frac{1}{\pi}\left[(\pi-t)\frac{\sin(nt)}{n}\right]_0^{2\pi} + \frac{1}{n\pi}\int_0^{2\pi}\sin(nt)\,\d t = 0
\]
\[
b_n = \frac{1}{\pi}\int_0^{2\pi}(\pi-t)\sin(nt)\,\d t = \frac{1}{\pi}\left[-(\pi-t)\frac{\cos(nt)}{n}\right]_0^{2\pi} - \frac{1}{n\pi}\int_0^{2\pi}\cos(nt)\,\d t
\]
\[
= \frac{1}{\pi n}\left[-(\pi-2\pi)(1) + (\pi-0)(1)\right] = \frac{1}{\pi n}(\pi + \pi)\cdot\frac{1}{1}
\]

Calcul direct : $b_n = \frac{1}{\pi}\left[-\frac{(\pi-t)\cos(nt)}{n}\right]_0^{2\pi} - \frac{1}{n\pi}\left[\frac{\sin(nt)}{n}\right]_0^{2\pi} = \frac{1}{n\pi}\left[\pi(-1) + \pi \right] + 0 = \dfrac{2}{n}$.

\textbf{2.} $f(t) = \displaystyle\sum_{n=1}^{+\infty}\frac{2}{n}\sin(nt) = 2\sin(t) + \sin(2t) + \frac{2}{3}\sin(3t)+\cdots$

\textbf{3.} En $t = 0$ : $f(0^+) = \pi$ et $f(2\pi^-) = \pi - 2\pi = -\pi$. La série converge vers $\dfrac{\pi + (-\pi)}{2} = 0$ (valeur moyenne). La discontinuité produit le phénomène de Gibbs : les sommes partielles oscillent d'environ 9\% au-dessus et en dessous de la discontinuité.

\textbf{4.} Parseval : $\dfrac{1}{2\pi}\int_0^{2\pi}(\pi-t)^2\,\d t = \dfrac{1}{2}\sum_{n=1}^\infty b_n^2$.

LHS : $\dfrac{1}{2\pi}\cdot\dfrac{[\text{-}(\pi\text{-}t)^3/3]_0^{2\pi}}{1} = \dfrac{1}{2\pi}\cdot\dfrac{\pi^3+\pi^3}{3} = \dfrac{\pi^2}{3}$.

RHS : $\dfrac{1}{2}\sum_{n=1}^\infty\dfrac{4}{n^2} = 2\sum_{n=1}^\infty\dfrac{1}{n^2}$.

Donc $\displaystyle\sum_{n=1}^{+\infty}\dfrac{1}{n^2} = \dfrac{\pi^2}{6}$.
\end{boitecorrige}

\subsection*{Exercice 11 - Résolution d'EDO par les séries de Fourier \niveau{\difficile}}

\begin{boiteexercice}
On cherche la solution $2\pi$-périodique de l'équation différentielle :
\[
y''(t) + \omega_0^2\,y(t) = f(t) = \frac{4}{\pi}\sum_{k=0}^{+\infty}\frac{\sin((2k+1)t)}{2k+1}
\]
où $f$ est le signal créneau (exercice 1) et $\omega_0 \notin \{1,3,5,\ldots\}$.
\begin{enumerate}
  \item En cherchant $y(t) = \sum_{n=1}^\infty B_n\sin(nt)$, substituer dans l'équation et identifier les $B_n$.
  \item Écrire la solution $y(t)$ sous forme de série.
  \item Que se passe-t-il si $\omega_0 = 1$ ? (phénomène de résonance)
\end{enumerate}
\end{boiteexercice}

\begin{boitecorrige}
\textbf{1.} On cherche $y(t) = \sum_{n=1}^\infty B_n\sin(nt)$. Alors $y'' = -\sum n^2 B_n\sin(nt)$. En substituant :
\[
\sum_{n=1}^\infty(-n^2+\omega_0^2)B_n\sin(nt) = \frac{4}{\pi}\sum_{k=0}^\infty\frac{\sin((2k+1)t)}{2k+1}
\]
Par identification des coefficients de $\sin(nt)$ :
\begin{itemize}[nosep]
  \item Pour $n = 2k+1$ (impair) : $(\omega_0^2 - n^2)B_n = \dfrac{4}{\pi n}$, donc $B_n = \dfrac{4}{\pi n(\omega_0^2 - n^2)}$.
  \item Pour $n$ pair : $B_n = 0$.
\end{itemize}

\textbf{2.} La solution est :
\[
\boxed{y(t) = \frac{4}{\pi}\sum_{k=0}^{+\infty}\frac{\sin((2k+1)t)}{(2k+1)(\omega_0^2-(2k+1)^2)}}
\]
Cette série converge si $\omega_0$ n'est pas un entier impair (pas de résonance).

\textbf{3.} Si $\omega_0 = 1$ : le terme $k=0$ donne un dénominateur nul : $\omega_0^2 - 1 = 0$. L'équation $y'' + y = \sin(t)$ est en résonance. La solution particulière n'est plus sinusoïdale mais de la forme $y_p = -\dfrac{t}{2}\cos(t)$, dont l'amplitude croît sans borne. C'est le phénomène de résonance.
\end{boitecorrige}


\subsection*{Exercice 12 - Signal carré et identité de Parseval \niveau{\moyen}}

\begin{boiteexercice}
Le signal créneau $f(t)$ de période $2\pi$, défini par $f(t) = +1$ pour $t \in ]0,\pi[$ et $f(t) = -1$ pour $t \in ]-\pi,0[$, admet la série de Fourier :
\[
f(t) = \frac{4}{\pi}\sum_{k=0}^{+\infty}\frac{\sin((2k+1)t)}{2k+1}
\]
\begin{enumerate}
  \item Rappeler l'identité de Parseval pour une série de Fourier : $\dfrac{1}{2\pi}\int_{-\pi}^{\pi}|f(t)|^2\,\d t = \dfrac{a_0^2}{4} + \dfrac{1}{2}\sum_{n=1}^{+\infty}(a_n^2+b_n^2)$.
  \item Calculer le membre de gauche (valeur efficace du signal créneau).
  \item Identifier les coefficients $b_{2k+1} = \dfrac{4}{\pi(2k+1)}$ et appliquer Parseval pour obtenir :
  \[
  \sum_{k=0}^{+\infty}\frac{1}{(2k+1)^2} = \frac{\pi^2}{8}
  \]
\end{enumerate}
\end{boiteexercice}

\begin{boitecorrige}
\textbf{1.} Pour une fonction $2\pi$-périodique de coefficients de Fourier $a_0, a_n, b_n$, l'identité de Parseval s'écrit :
\[
\frac{1}{2\pi}\int_{-\pi}^{\pi}|f(t)|^2\,\d t = \frac{a_0^2}{4} + \frac{1}{2}\sum_{n=1}^{+\infty}(a_n^2+b_n^2)
\]

\textbf{2.} Pour le signal créneau, $|f(t)|^2 = 1$ sur tout $]-\pi,\pi[$ (puisque $f(t) = \pm 1$). Donc :
\[
\frac{1}{2\pi}\int_{-\pi}^{\pi} 1\,\d t = \frac{1}{2\pi}\times 2\pi = 1
\]

\textbf{3.} Le signal créneau est impair, donc $a_0 = 0$ et tous les $a_n = 0$. Seuls les $b_{2k+1} = \dfrac{4}{\pi(2k+1)}$ sont non nuls (coefficients impairs). Parseval donne :
\[
1 = \frac{1}{2}\sum_{k=0}^{+\infty}b_{2k+1}^2 = \frac{1}{2}\sum_{k=0}^{+\infty}\left(\frac{4}{\pi(2k+1)}\right)^2 = \frac{1}{2}\cdot\frac{16}{\pi^2}\sum_{k=0}^{+\infty}\frac{1}{(2k+1)^2}
\]
\[
\boxed{\sum_{k=0}^{+\infty}\frac{1}{(2k+1)^2} = \frac{\pi^2}{8}}
\]
\end{boitecorrige}

\subsection*{Exercice 13 - Transformée de Fourier d'un signal exponentiel \niveau{\moyen}}

\begin{boiteexercice}
On considère le signal causal $f(t) = e^{-at}\,u(t)$ où $a > 0$ et $u(t)$ est l'échelon de Heaviside ($u(t) = 1$ pour $t \geq 0$, $u(t) = 0$ pour $t < 0$).
\begin{enumerate}
  \item Rappeler la définition de la transformée de Fourier $F(\omega) = \displaystyle\int_{-\infty}^{+\infty} f(t)\,e^{-j\omega t}\,\d t$.
  \item Calculer $F(\omega)$ et montrer que $F(\omega) = \dfrac{1}{a + j\omega}$.
  \item Calculer le module $|F(\omega)| = \dfrac{1}{\sqrt{a^2+\omega^2}}$ et la phase $\arg(F(\omega)) = -\arctan(\omega/a)$.
  \item Pour $a = 1$, calculer $|F(0)|$, $|F(a)|$ et $|F(10a)|$.
\end{enumerate}
\end{boiteexercice}

\begin{boitecorrige}
\textbf{1.} La transformée de Fourier est définie par :
\[
F(\omega) = \int_{-\infty}^{+\infty} f(t)\,e^{-j\omega t}\,\d t
\]

\textbf{2.} Comme $f(t) = 0$ pour $t < 0$ et $f(t) = e^{-at}$ pour $t \geq 0$ :
\[
F(\omega) = \int_0^{+\infty} e^{-at}\,e^{-j\omega t}\,\d t = \int_0^{+\infty} e^{-(a+j\omega)t}\,\d t = \left[\frac{e^{-(a+j\omega)t}}{-(a+j\omega)}\right]_0^{+\infty}
\]
Comme $a > 0$, l'exponentielle converge vers $0$ à l'infini :
\[
\boxed{F(\omega) = \frac{1}{a+j\omega}}
\]

\textbf{3.} Module et phase :
\[
|F(\omega)| = \frac{1}{|a+j\omega|} = \frac{1}{\sqrt{a^2+\omega^2}}
\]
\[
\arg(F(\omega)) = \arg(1) - \arg(a+j\omega) = 0 - \arctan\!\left(\frac{\omega}{a}\right) = -\arctan\!\left(\frac{\omega}{a}\right)
\]

\textbf{4.} Pour $a = 1$ :
\begin{align*}
|F(0)|  &= 1,\\
|F(1)|  &= \frac{1}{\sqrt{2}} \approx 0{,}707,\\
|F(10)| &\approx 0{,}0995
\end{align*}
La fréquence de coupure est $\omega_c = a$
(le module est divisé par $\sqrt{2}$ par rapport à la valeur en $\omega = 0$).
\end{boitecorrige}

\subsection*{Exercice 14 - Convolution par transformée de Fourier \niveau{\difficile}}

\begin{boiteexercice}
On donne deux signaux causaux : $x(t) = e^{-2t}\,u(t)$ et $h(t) = e^{-t}\,u(t)$.
\begin{enumerate}
  \item Calculer les transformées de Fourier $X(\omega)$ et $H(\omega)$.
  \item En utilisant la propriété de convolution ($\mathcal{F}\{x*h\} = X\cdot H$), exprimer $Y(\omega) = X(\omega)\cdot H(\omega)$.
  \item Décomposer $Y(\omega)$ en éléments simples.
  \item En utilisant le résultat de l'exercice précédent, calculer la transformée de Fourier inverse $y(t) = (x*h)(t)$.
\end{enumerate}
\end{boiteexercice}

\begin{boitecorrige}
\textbf{1.} Par l'exercice précédent (formule $e^{-at}u(t) \xrightarrow{\mathcal{F}} 1/(a+j\omega)$) :
\[
X(\omega) = \frac{1}{2+j\omega}, \qquad H(\omega) = \frac{1}{1+j\omega}
\]

\textbf{2.} Produit :
\[
Y(\omega) = X(\omega)\cdot H(\omega) = \frac{1}{(2+j\omega)(1+j\omega)}
\]

\textbf{3.} Décomposition en éléments simples (en posant $p = j\omega$) :
\[
\frac{1}{(1+p)(2+p)} = \frac{A}{1+p} + \frac{B}{2+p}
\]
Calcul des résidus : $A = \displaystyle\lim_{p\to-1}(1+p)\cdot\frac{1}{(1+p)(2+p)} = \frac{1}{1} = 1$ et $B = \displaystyle\lim_{p\to-2}\frac{1}{1+p} = \frac{1}{-1} = -1$.

Donc :
\[
Y(\omega) = \frac{1}{1+j\omega} - \frac{1}{2+j\omega}
\]

\textbf{4.} Par linéarité et en utilisant le résultat $1/(a+j\omega) \xrightarrow{\mathcal{F}^{-1}} e^{-at}u(t)$ :
\[
\boxed{y(t) = e^{-t}u(t) - e^{-2t}u(t) = \left(e^{-t} - e^{-2t}\right)u(t)}
\]
Vérification : $y(0) = 1 - 1 = 0$, $y(t) > 0$ pour $t > 0$ (car $e^{-t} > e^{-2t}$), $y(t) \to 0$ pour $t \to +\infty$.
\end{boitecorrige}

\ifdefined\inMainDoc\else\end{document}\fi
